\begin{proof}[Proof of Theorem~\ref{thm:localizedcase}]
We must show that $L(\Gamma, E)|_{F}$ is non-zero in degree $|F_1|$. 
Recall that $L(\Gamma, E)|_{F}$ is isomorphic to $I|_{F}/J|_{F}$, where $I|_{F}$ and $J|_{F}$ are described in \eqref{eq:Ires} and \eqref{eq:Jres} respectively. First we handle the cases when $\vert F_1 \vert \le 1$. If $F_1 = \emptyset$, then $E$ is interior and $x^{\emptyset} = 1$, but $J|_{F}$ is a proper ideal as it is generated by elements of positive degree, so $x^{F_1} \not \in J|_{F}$. If $F_1 = \{v\}$, then we assume that $E$ is not an interior face. Then $J|_{F}$ is generated by elements of degree at least $2$, so $x^{F_1} \not \in J|_{F}$.

Suppose $\vert F_1 \vert = 2$. We assume that there are no vertices $v$ with $\{v\} \sqcup E$ interior and $E$ is not interior. If $\sigma(E)$ has codimension $1$, then both $F_1$ and $F_2$ must have a vertex $v$ with $\{v\} \sqcup E$ interior. Then by Lemma~\ref{lem:cd1case}, we see that $\dim L(\Gamma, E)|_{F} \ge 1$. Hence we may assume that $\sigma(E)$ has codimension at least $2$.

Let $F_1 = \{u, t\}$ and assume that $L(\Gamma, E)|_{F}$ has no non-zero elements in degree $2$. Consider the connected component of the internal edge graph containing $F_1$. By Lemma~\ref{lem:intedge}, we may assume that $\sigma(\{u\} \sqcup E) = \{v_1\}^c$. Note that $v_1 \in \sigma(t)$. There is a vertex $t' \in F_2$ such that $v_1 \in \sigma(t')$, so $\{u, t'\} \sqcup E$ is interior. Therefore either $\{t\} \sqcup E$ or $\{t'\} \sqcup E$ has carrier codimension $1$. 

If $\sigma(\{t\} \sqcup E) = \{v_2\}^c$, then there is a vertex $u' \in F_2$ such that $v_2 \in \sigma(u')$. First assume $u'$ and $t'$ are distinct. Since at least one of $\{u'\} \sqcup E$ and $\{t'\} \sqcup E$ has carrier codimension $1$, it follows that $\{ u', t'\} \sqcup E$ is interior. Then $\{u, t, u', t'\}$ forms a $4$-cycle, contradicting Lemma~\ref{lem:nofourcycle}.

If $u' = t'$, then the internal edge graph contains a cycle and hence every vertex $w$ in it (including $t$) has $\{w\} \sqcup E$ of carrier codimension $1$. As $F_2$ is interior and $\{u'\} \sqcup E$ has carrier codimension $1$, there is a vertex $w \in F_2$ such that $\{u', w\} \sqcup E$ is interior. But then either $\{u, w\} \sqcup E$ or $\{t, w\} \sqcup E$ is interior, contradicting the uniqueness of the cycle in Lemma~\ref{lem:intedge}.

If $\{t\} \sqcup E$ does not have carrier codimension $1$, then we may assume that $\sigma(\{t'\} \sqcup E) = \{v_2\}^c$. Choose a vertex $u' \in F_2$ with $v_2 \in \sigma(u')$. Then $\{t', u'\} \sqcup E$ is interior, so $\{u'\} \sqcup E$ has carrier codimension $1$. If $v_1 \in \sigma(u')$, then $\{u, u'\} \sqcup E$ is interior. If $v_1 \not \in \sigma(u')$, then $\{t, u'\} \sqcup E$ is interior. In either case, there is a cycle and a vertex $v$ with $\{v\} \sqcup E$ of carrier codimension more than $1$ in the internal edge graph, contradicting Lemma~\ref{lem:intedge}. 
\end{proof}
